\slcol{श्रीभगवानुवाच —\\
\Index{इदं तु ते गुह्यतमं} प्रवक्ष्याम्यनसूयवे । \\
ज्ञानं विज्ञानसहितं यज्ज्ञात्वा मोक्ष्यसेऽशुभात् ॥ ९.१ ॥ }
\translate{śrībhagavānuvāca —\\
idaṁ tu tē guhyatamaṁ pravakṣyāmyanasūyavē | \\
jñānaṁ vijñānasahitaṁ yajjñātvā mōkṣyasē:'śubhāt || 9.1 || }
\cquote{Sri Bhagavan said: I shall now declare to you, who are free from envy, this most secret knowledge, together with wisdom (its realisation). By knowing this, you shall be liberated from all that is inauspicious.}
\slcol{\Index{राजविद्या राजगुह्यं} पवित्रमिदमुत्तमम् । \\
प्रत्यक्षावगमं धर्म्यं सुसुखं कर्तुमव्ययम् ॥ ९.२ ॥ }
\translate{rājavidyā rājaguhyaṁ pavitramidamuttamam | \\
pratyakṣāvagamaṁ dharmyaṁ \\susukhaṁ kartumavyayam || 9.2 || }
\cquote{This royal knowledge (the king of sciences), the most profound of secrets, is supremely pure, directly realisable, righteous, easy to practise, and imperishable.}
\newpage
\begin{mananam}{\mananamfont{Mananam Śloka - 1, 2}}
\normalfont\mananamtext
Why is spiritual knowledge spoken of as hidden or secret, even though it is widely available? What qualities or traits might make someone ready to receive spiritual teachings? How does the progression to advanced levels of spiritual knowledge compare to a student progressing from high school to college? What are the similarities and differences in these paradigms? 
How might theoretical knowledge differ from spiritual wisdom? What signs or qualities might indicate that someone has attained true wisdom or the knowledge of spiritual realization? How might the attitudes and behaviours of those who are wise differ from those who only hold theoretical knowledge?
\end{mananam}
\sreflect
\begin{inspiration}{\inspirationTitleFont{Inspiration}}
\normalfont\mananamtext The sacred knowledge of spirituality is accessible only to those with a ``royal mind''—one that is pure, surrendered, and disciplined. To attain this wisdom, one must cultivate a deep attunement with the \emph{guru} and the tradition.
In a sense, anything unknown to us remains a secret. Anything we haven't sought to discover or have not yet found a teacher to reveal can also be considered hidden Yet, in the boundlessness of God’s mind, every mystery is waiting to be uncovered by those who are truly prepared and earnest in their search.
\end{inspiration}
\newpage
\slcol{\Index{अश्रद्दधानाः पुरुषा} धर्मस्यास्य परन्तप । \\
अप्राप्य मां निवर्तन्ते मृत्युसंसारवर्त्मनि ॥ ९.३ ॥ }
\translate{aśraddadhānāḥ puruṣā dharmasyāsya parantapa | \\
aprāpya māṁ nivartantē mr̥tyusaṁsāravartmani || 9.3 || }
\cquote{O Parantapa (scorcher of foes, Arjuna), those who lack faith in this \emph{dharma} do not attain Me. They return to the path of death and rebirth in the cycle of \emph{saṁsāra}.}
\slcol{\Index{मया ततमिदं सर्वं} जगतदव्यक्तमूर्तिना । \\
मत्स्थानि सर्वभूतानि न चाहं तेष्ववस्थितः ॥ ९.४ ॥ }
\translate{mayā tatamidaṁ sarvaṁ jagatadavyaktamūrtinā | \\
matsthāni sarvabhūtāni na cāhaṁ tēṣvavasthitaḥ || 9.4 || }
\cquote{By Me, in My unmanifest form, this entire universe is pervaded. All beings dwell in Me, yet I do not dwell in them (am not contained in them).}
\slcol{\Index{न च मत्स्थानि भूतानि} पश्य मे योगमैश्वरम् । \\
भूतभृन्न च भूतस्थो ममात्मा भूतभावनः ॥ ९.५ ॥ }
\translate{na ca matsthāni bhūtāni paśya mē yōgamaiśvaram | \\
bhūtabhr̥nna ca bhūtasthō mamātmā bhūtabhāvanaḥ || 9.5 || }
\cquote{Nor do beings truly dwell in Me. Behold My divine \emph{yoga}! Though I sustain all beings, I am my Self not contained within them. And yet, I am the source and nurturer of all beings.}
\slcol{\Index{यथाकाशस्थितो नित्यं} वायुः सर्वत्रगो महान् । \\
तथा सर्वाणि भूतानि मत्स्थानीत्युपधारय ॥ ९.६ ॥ }
\translate{yathākāśasthitō nityaṁ vāyuḥ sarvatragō mahān | \\
tathā sarvāṇi bhūtāni matsthānītyupadhāraya || 9.6 || }
\cquote{Just as the mighty wind, moving everywhere, always rests in space (\emph{ākāśa}), so too, understand that all beings dwell in Me.}

\newpage
\begin{mananam}{\mananamfont{Mananam Śloka - 4, 5, 6}}
\normalfont\mananamtext What hidden spiritual and scientific principles can we learn from observing space? Should we dismiss space as non-existent simply because it isn’t as tangible as earth, water, or fire? Could any of nature’s elements exist without space? Would wind move, fire rage, or water flow if space were non-existent? Yet, isn’t space distinct and transcendent, beyond all other elements?
In the same way, might it be possible to conceive a divine entity that sustains all existence while remaining distinct and transcendent?
\end{mananam}
\sreflect
\begin{inspiration}{\inspirationTitleFont{Inspiration}}
\normalfont\mananamtext The forces that uphold matter and are the support of all nature are not readily available to scientists despite their finest instruments and technology. Yet, the science behind the creation of this whole universe is quite easily accessible to those who have a subtle sense of observing, and who use the pointers left behind by the ancient sages who had fine inner instruments with which they could launch into a close observation of nature and its source.
\end{inspiration}
\newpage

\slcol{\Index{सर्वभूतानि कौन्तेय} प्रकृतिं यान्ति मामिकाम् । \\
कल्पक्षये पुनस्तानि कल्पादौ विसृजाम्यहम् ॥ ९.७ ॥ }
\translate{sarvabhūtāni kauntēya prakr̥tiṁ yānti māmikām | \\
kalpakṣayē punastāni kalpādau visr̥jāmyaham || 9.7 || }
\cquote{O Kaunteya (son of Kunti, Arjuna), at the end of a \emph{kalpa} (cosmic cycle) all beings enter into My \emph{prakṛti}; and at the beginning of the next \emph{kalpa}, I send them forth again.}
\slcol{\Index{प्रकृतिं स्वामवष्टभ्य} विसृजामि पुनः पुनः । \\
भूतग्राममिमं कृत्स्नमवशं प्रकृतेर्वशात् ॥ ९.८ ॥ }
\translate{prakr̥tiṁ svāmavaṣṭabhya visr̥jāmi punaḥ punaḥ | \\
bhūtagrāmamimaṁ kr̥tsnamavaśaṁ prakr̥tērvaśāt || 9.8 || }
\cquote{Controlling My own \emph{prakṛti}, I send forth again and again this entire multitude of beings, who are helplessly bound by the power of \emph{prakṛti} (through \emph{guṇas} and \emph{karma}).}
\slcol{\Index{न च मां तानि कर्माणि} निबध्नन्ति धनञ्जय । \\
उदासीनवदासीनमसक्तं तेषु कर्मसु ॥ ९.९ ॥ }
\translate{na ca māṁ tāni karmāṇi nibadhnanti dhanañjaya | \\
udāsīnavadāsīnamasaktaṁ tēṣu karmasu || 9.9 || }
\cquote{O Dhanañjaya (the winner of wealth, Arjuna), those actions do not bind Me. Remaining as if indifferent, unattached to those actions, I sit apart from them.}
\slcol{\Index{मयाध्यक्षेण प्रकृतिः} सूयते सचराचरम् । \\
हेतुनानेन कौन्तेय जगद्विपरिवर्तते ॥ ९.१० ॥ }
\translate{mayādhyakṣēṇa prakr̥tiḥ sūyatē sacarācaram | \\
hētunānēna kauntēya jagadviparivartatē || 9.10 || }
\cquote{Under My supervision, \emph{prakṛti} brings forth the moving and the unmoving; by this cause, O Kaunteya, the world revolves.}
\slcol{\Index{अवजानन्ति मां मूढा} मानुषीं तनुमाश्रितम् । \\
परं भावमजानन्तो मम भूतमहेश्वरम् ॥ ९.११ ॥ }
\translate{avajānanti māṁ mūḍhā mānuṣīṁ tanumāśritam | \\
paraṁ bhāvamajānantō mama bhūtamahēśvaram || 9.11 || }
\cquote{Deluded men disregard Me when I dwell in a human form, not knowing My higher nature as the great Lord of all beings.}
\slcol{\Index{मोघाशा मोघकर्माणो} मोघज्ञाना विचेतसः । \\
राक्षसीमासुरीं चैव प्रकृतिं मोहिनीं श्रिताः ॥ ९.१२ ॥ }
\translate{mōghāśā mōghakarmāṇō mōghajñānā vicētasaḥ | \\
rākṣasīmāsurīṁ caiva prakr̥tiṁ mōhinīṁ śritāḥ || 9.12 || }
\cquote{Vain are their hopes, vain are their actions, vain are their knowledge; deluded in mind, they take refuge in the bewildering \emph{prakṛti} nature of \emph{rākṣasas} and \emph{asuras}.}
\slcol{\Index{महात्मानस्तु मां पार्थ} दैवीं प्रकृतिमाश्रिताः । \\
भजन्त्यनन्यमनसो ज्ञात्वा भूतादिमव्ययम् ॥ ९.१३ ॥ }
\translate{mahātmānastu māṁ pārtha daivīṁ prakr̥timāśritāḥ | \\
bhajantyananyamanasō jñātvā bhūtādimavyayam || 9.13 || }
\cquote{But the \emph{mahatmas} (great souls), O Pārtha, taking refuge in the divine \emph{prakṛti}, worship Me with single-minded devotion, knowing Me as the imperishable origin of all beings.}
\slcol{\Index{सततं कीर्तयन्तो मां} यतन्तश्च दृढव्रताः । \\
नमस्यन्तश्च मां भक्त्या नित्ययुक्ता उपासते ॥ ९.१४ ॥ }
\translate{satataṁ kīrtayantō māṁ yatantaśca dr̥ḍhavratāḥ | \\
namasyantaśca māṁ bhaktyā nityayuktā upāsatē || 9.14 || }
\cquote{Always chanting My glories, striving with firm vows, bowing down to Me with devotion, the ever-steadfast worship Me.}
\slcol{\Index{ज्ञानयज्ञेन चाप्यन्ये} यजन्तो मामुपासते । \\
एकत्वेन पृथक्त्वेन बहुधा विश्वतोमुखम् ॥ ९.१५ ॥ }
\translate{jñānayajñēna cāpyanyē yajantō māmupāsatē | \\
ēkatvēna pr̥thaktvēna bahudhā viśvatōmukham || 9.15 || }
\cquote{Others, too, worship Me through the knowledge-sacrifice (\emph{jñāna-yajña}); by seeing Me as one, as separate, and in many forms—as the Universal-faced (with faces everywhere).}

\newpage
\begin{mananam}{\mananamfont{Mananam Śloka - 12, 13}}
\normalfont\mananamtext Where do my desires and aspirations come from? To what extent do my actions stem from my needs and wants? Do I pursue things based on my own direct knowledge, or am I influenced by what others claim is worth seeking? Do I pause to consider whether what I pursue will truly bring me happiness? What is the source of my beliefs—my culture, my upbringing, my experiences? How can I trace back to the origins of my understanding, my beliefs, and their formation? Can ignorance—of the higher reality of my own being and of God—truly be bliss?
\end{mananam}
\sreflect
\begin{inspiration}{\inspirationTitleFont{Inspiration}}
\normalfont\mananamtext What leads one person to atheism, another to theism, and others to become seekers? It is the dimension of their knowledge and understanding. Our pursuits and actions in life mirror the scope of our awareness. Without clear spiritual insight, we often focus solely on personal well-being, which reflects a state of spiritual ignorance.
Staying content in ignorance may feel comfortable, but it is not lasting. Chasing fleeting moments of happiness, knowing that they provide only temporary satisfaction, means missing out on the potential for the happiness that endures and which comes from actions guided by the wisdom of a higher reality.
\end{inspiration}
\newpage

\slcol{\Index{अहं क्रतुरहं यज्ञः} स्वधाहमहमौषधम् । \\
मन्त्रोऽहमहमेवाज्यमहमग्निरहं हुतम् ॥ ९.१६ ॥ }
\translate{ahaṁ kraturahaṁ yajñaḥ svadhāhamahamauṣadham | \\
mantrō:'hamahamēvājyamahamagnirahaṁ hutam || 9.16 || }
\cquote{I am the \emph{kratu} (\emph{vedic} ritual), I am the \emph{yajña} (sacrifice), I am the \emph{svadhā} (offering to the ancestors), I am the medicinal herb, I am the \emph{mantra}, I am the clarified butter, I am the fire, and I am the oblation.}
\slcol{\Index{पिताहमस्य जगतो} माता धाता पितामहः । \\
वेद्यं पवित्रमोङ्कार ऋक्साम यजुरेव च ॥ ९.१७ ॥ }
\translate{pitāhamasya jagatō mātā dhātā pitāmahaḥ | \\
vēdyaṁ pavitramōṅkāra r̥ksāma yajurēva ca || 9.17 || }
\cquote{I am the father of this world, the mother, the sustainer, and the grandsire. I am the object of knowledge, the purifier, the \emph{oṁkāra}, and also the \emph{Ṛg, Sāma}, and \emph{Yajur Vedas}.}
\slcol{\Index{गतिर्भर्ता प्रभुः साक्षी} निवासः शरणं सुहृत् । \\
प्रभवः प्रलयः स्थानं निधानं बीजमव्ययम् ॥ ९.१८ ॥ }
\translate{gatirbhartā prabhuḥ sākṣī nivāsaḥ śaraṇaṁ suhr̥t | \\
prabhavaḥ pralayaḥ sthānaṁ nidhānaṁ bījamavyayam || 9.18 ||}
\cquote{I am the goal, the sustainer, the Lord, the witness, the abode, the refuge, the friend. I am the origin, the dissolution, the foundation, the treasure, and the imperishable seed.}
\slcol{\Index{तपाम्यहमहं वर्षं} निगृह्णाम्युत्सृजामि च । \\
अमृतं चैव मृत्युश्च सदसच्चाहमर्जुन ॥ ९.१९ ॥ }
\translate{tapāmyahamahaṁ varṣaṁ nigr̥hṇāmyutsr̥jāmi ca | \\
amr̥taṁ caiva mr̥tyuśca sadasaccāhamarjuna || 9.19 || }
\cquote{I am the source of heat; I bring forth rain and also withhold it. I am immortality and also death, and existence and non-existence, O Arjuna.}

\newpage
\begin{mananam}{\mananamfont{Mananam Śloka - 18, 19}}
\normalfont\mananamtext What is the nature of the ultimate source of all creation? Could it also be the final destination for all beings? Could the entity that initiated creation also be responsible for its dissolution, as some scriptures suggest? 
How might that entity exert influence over me and other beings? Could it be the controlling force behind my birth, life, and death? What value is there in seeking answers to these questions? Is it better to let others uncover these truths and reveal them to me? And, in matters like these, can the answers that others provide me be truly satisfying? 
\end{mananam}
\sreflect
\begin{inspiration}{\inspirationTitleFont{Inspiration}}
\normalfont\mananamtext From time immemorial, humanity has sought to unravel the origin of all creation. Theories abound—from belief in a single controlling entity to the idea that such a single creator doesn’t exist—captivating scientists and philosophers alike. Through our natural tendency for cause-hunting, it seems logical to infer that there must be an ultimate source from which even the Big Bang emerged. Yet, how can we be certain that this ultimate cause itself has no prior cause? To reconcile these possibilities, the Lord encompasses both ideas: that of a personal God, and that of and an impersonal Reality beyond causation—a transcendental truth of which the universe is merely a shadow. 
\end{inspiration}
\newpage

\slcol{\Index{त्रैविद्या मां सोमपाः} पूतपापा यज्ञैरिष्ट्वा स्वर्गतिं प्रार्थयन्ते ।\\
ते पुण्यमासाद्य सुरेन्द्रलोकमश्नन्ति दिव्यान्दिवि देवभोगान् ॥ ९.२० ॥ }
\translate{traividyā māṁ sōmapāḥ pūtapāpā\\ \hspace*{0.5cm}yajñairiṣṭvā svargatiṁ prārthayantē |\\
tē puṇyamāsādya surēndralōkamaśnanti\\ \hspace*{0.5cm}divyāndivi dēvabhōgān || 9.20 ||}
\cquote{Those who are versed in the three \emph{Vedas}, who drink the \emph{soma} and are purified of sin, worship Me through sacrifices and seek heaven. Attaining the world of Indra (the Lord of demi-gods), they enjoy the divine pleasures of the gods in heaven.}
\slcol{\Index{ते तं भुक्त्वा स्वर्गलोकं} विशालं क्षीणे पुण्ये मर्त्यलोकं विशन्ति । \\
एवं त्रयीधर्ममनुप्रपन्ना गतागतं कामकामा लभन्ते ॥ ९.२१ ॥}
\translate{tē taṁ bhuktvā svargalōkaṁ viśālaṁ \\
\hspace*{0.5cm}kṣīṇē puṇyē martyalōkaṁ viśanti | \\
ēvaṁ trayīdharmamanuprapannā \\
\hspace*{0.5cm}gatāgataṁ kāmakāmā labhantē || 9.21 || }
\cquote{They, having enjoyed the vast heavenly world, when their merit is exhausted, return to the mortal realm. Thus, those who follow the threefold \emph{Vedic} path, desiring pleasures, attain only repeated coming and going (of \emph{saṁsāra}).}
\slcol{\Index{अनन्याश्चिन्तयन्तो मां} ये जनाः पर्युपासते । \\
तेषां नित्याभियुक्तानां योगक्षेमं वहाम्यहम् ॥ ९.२२ ॥ }
\translate{ananyāścintayantō māṁ yē janāḥ paryupāsatē |\\ 
tēṣāṁ nityābhiyuktānāṁ yōgakṣēmaṁ vahāmyaham || 9.22 || }
\cquote{Those who, thinking of none other, worship Me alone with undivided devotion—to them, ever steadfast, I provide what they lack (\emph{yoga}) and preserve what they already have (\emph{kṣema}).}

\newpage
\begin{mananam}{\mananamfont {Mananam Śloka - 22}}
\normalfont\mananamtext Which of life’s two pursuits engages my mind more—acquiring what I do not yet have, or preserving what I already possess? What specific aspects fall into each of these categories for me? How much control do I truly have over elements such as health, wealth, social status, and success? Can I place my trust in a higher justice system or divine order to provide for my needs and safeguard me, while I make an earnest effort to live righteously?
\end{mananam}
\sreflect
\begin{inspiration}{\inspirationTitleFont{Inspiration}}
\normalfont\mananamtext We all have a choice: to immerse ourselves in chasing our needs and desires, or to place our focus on the Divine. Practical-minded people might suggest prioritising worldly concerns, as these seem more immediate and tangible. However, the Lord assures us that when we fulfil our duties and dedicate time daily for genuine contemplation and devotion, He will provide for our needs and protect our worldly gains. For those who have faith, trusting in the Divine brings both spiritual and material rewards.
\end{inspiration}
\newpage

\slcol{\Index{येऽप्यन्यदेवताभक्ता} यजन्ते श्रद्धयान्विताः । \\
तेऽपि मामेव कौन्तेय यजन्त्यविधिपूर्वकम् ॥ ९.२३ ॥ }
\translate{yē:'pyanyadēvatābhaktā yajantē śraddhayānvitāḥ | \\
tē:'pi māmēva kauntēya yajantyavidhipūrvakam || 9.23 || }
\cquote{Even those devotees who, with faith, worship other deities, O Kaunteya, they too worship Me alone—but in a manner not according to the proper method.}
\slcol{\Index{अहं हि सर्वयज्ञानां} भोक्ता च प्रभुरेव च । \\
न तु मामभिजानन्ति तत्त्वेनातश्च्यवन्ति ते ॥ ९.२४ ॥ }
\translate{ahaṁ hi sarvayajñānāṁ bhōktā ca prabhurēva ca | \\
na tu māmabhijānanti tattvēnātaścyavanti tē || 9.24 || }
\cquote{For I am indeed the enjoyer and the Lord of all sacrifices. But they do not know Me in truth; therefore, they fall away (do not gain liberation).}
\slcol{\Index{यान्ति देवव्रता} देवान्पितॄन्यान्ति पितृव्रताः । \\
भूतानि यान्ति भूतेज्या यान्ति मद्याजिनोऽपि माम् ॥ ९.२५ ॥ }
\translate{yānti dēvavratā dēvānpitr̥̄nyānti pitr̥vratāḥ | \\
bhūtāni yānti bhūtējyā yānti madyājinō:'pi mām || 9.25 || }
\cquote{Worshippers of the gods (\emph{devas}) go to the gods; worshippers of the ancestors (\emph{pitṝs}) go to the ancestors; worshippers of spirits (\emph{bhūtās}) go to the spirits; but those who worship Me come to Me.}
\slcol{\Index{पत्रं पुष्पं फलं तोयं} यो मे भक्त्या प्रयच्छति । \\
तदहं भक्त्युपहृतमश्नामि प्रयतात्मनः ॥ ९.२६ ॥ }
\translate{patraṁ puṣpaṁ phalaṁ tōyaṁ yō mē bhaktyā prayacchati | \\
tadahaṁ bhaktyupahr̥tamaśnāmi prayatātmanaḥ || 9.26 || }
\cquote{Whoever offers Me with devotion a leaf, a flower, a fruit, or water—I accept that offering, made with devotion by one of pure heart.}

\newpage
\begin{mananam}{\mananamfont {Mananam Śloka - 25, 26}}
\normalfont\mananamtext What motivates you to perform good deeds or rituals? Do you do them simply because you’ve been told they will bring specific benefits? Do you possess enough understanding to pray or offer your devotion to the Divine with complete trust and sincerity? 
Have you ever felt gratitude for a kind gesture from a loved one? In the same way, could you offer something as simple as a flower or fruit to the Lord, filled with pure devotion? When you do this, doesn’t it leave an uplifting, transformative effect on your state of mind?
\end{mananam}
\sreflect
\begin{inspiration}{\inspirationTitleFont{Inspiration}}
\normalfont\mananamtext Clear understanding of physical, verbal, and mental acts of worship and their purpose is superior to merely going through their ritualistic steps. When our devotion is directed toward lesser beings or deities who can fulfil our desires, we become indebted to them. However, if our devotion is toward the one Supreme God, the essence of our being, we are blessed with spiritual realisation and liberation.
When we purify our hearts of worldly desires and offer anything to God or to holy persons, we are rewarded manifold. Even if we offer nothing material, entering deeply into the Divine presence with a pure heart is a powerful remedy for psychological distress, as it brings an immediate, transformative sense of peace and joy.
\end{inspiration}
\newpage

\slcol{\Index{यत्करोषि यदश्नासि} यज्जुहोषि ददासि यत् । \\
यत्तपस्यसि कौन्तेय तत्कुरुष्व मदर्पणम् ॥ ९.२७ ॥ }
\translate{yatkarōṣi yadaśnāsi yajjuhōṣi dadāsi yat | \\
yattapasyasi kauntēya tatkuruṣva madarpaṇam || 9.27 || }
\cquote{Whatever you do, whatever you eat, whatever you offer in sacrifice, whatever you give, whatever austerity you practise—O Kaunteya, do that as an offering to Me.}
\slcol{\Index{शुभाशुभफलैरेवं} मोक्ष्यसे कर्मबन्धनैः । \\
संन्यासयोगयुक्तात्मा विमुक्तो मामुपैष्यसि ॥ ९.२८ ॥ }
\translate{śubhāśubhaphalairēvaṁ mōkṣyasē karmabandhanaiḥ | \\
saṁnyāsayōgayuktātmā vimuktō māmupaiṣyasi || 9.28 || }
\cquote{Thus, you shall be freed from the bonds of action, with their auspicious and inauspicious fruits. With the mind united in the \emph{yoga} of renunciation, liberated, you shall come to Me.}
\slcol{\Index{समोऽहं सर्वभूतेषु} न मे द्वेष्योऽस्ति न प्रियः । \\
ये भजन्ति तु मां भक्त्या मयि ते तेषु चाप्यहम् ॥ ९.२९ ॥ }
\translate{samō:'haṁ sarvabhūtēṣu na mē dvēṣyō:'sti na priyaḥ | \\
yē bhajanti tu māṁ bhaktyā mayi tē tēṣu cāpyaham || 9.29 || }
\cquote{I am equal to all beings; none is hateful to Me, nor is anyone dearer than another. Yet those who worship Me with devotion—they are in Me, and I am in them.}
\slcol{\Index{अपि चेत्सुदुराचारो} भजते मामनन्यभाक् । \\
साधुरेव स मन्तव्यः सम्यग्व्यवसितो हि सः ॥ ९.३० ॥ }
\translate{api cētsudurācārō bhajatē māmananyabhāk | \\
sādhurēva sa mantavyaḥ samyagvyavasitō hi saḥ || 9.30 || }
\cquote{Even if a person of very sinful conduct worships Me with undivided devotion (or with devotion to none else), he should be regarded as righteous, for such a one has rightly resolved.}

\newpage
\begin{mananam}{\mananamfont {Mananam Śloka - 29, 30}}
\normalfont\mananamtext Could the Lord of the universe show favouritism towards certain individuals? What determines who receives grace and blessings? Is there a difference between divine grace and the results of \emph{karma}? What is the significance of being devoted to God and surrendering our actions to Him? How might this mindset shape our actions and affect the rewards we receive?
Is there no hope for the blessings of spiritual growth for those new to the spiritual path, or for those who find it difficult to maintain faith? Are individuals who lack faith inferior in their essential nature to those naturally gifted with it?
\end{mananam}
\sreflect
\begin{inspiration}{\inspirationTitleFont{Inspiration}}
\normalfont\mananamtext If the Lord were to show bias or favouritism, He would be more human than divine. While partiality is common among businessmen and some political leaders, it is not fitting for the Supreme Lord of the Universe. However, those who adopt the right attitude, surrender to the Divine, and follow the path of righteousness are rightly rewarded with special grace and blessings.
Even those who stray from the righteous path can receive divine grace as soon as they change their ways, much like a mischievous child who is immediately forgiven and embraced by a loving mother when she corrects his behaviour.
\end{inspiration}
\newpage

\slcol{\Index{क्षिप्रं भवति धर्मात्मा} शश्वच्छान्तिं निगच्छति । \\
कौन्तेय प्रतिजानीहि न मे भक्तः प्रणश्यति ॥ ९.३१ ॥ }
\translate{kṣipraṁ bhavati dharmātmā 
śaśvacchāntiṁ nigacchati | \\
kauntēya pratijānīhi 
na mē bhaktaḥ praṇaśyati || 9.31 || }
\cquote{Such a one quickly becomes righteous (undergoing self-transformation), and attains eternal peace. O Kaunteya, know with certainty: My devotee never perishes.}
\slcol{\Index{मां हि पार्थ व्यपाश्रित्य} येऽपि स्युः पापयोनयः । \\
स्त्रियो वैश्यास्तथा शूद्रास्तेऽपि यान्ति परां गतिम् ॥ ९.३२ ॥ }
\translate{māṁ hi pārtha vyapāśritya yē:'pi syuḥ pāpayōnayaḥ | \\
striyō vaiśyāstathā śūdrāstē:'pi yānti parāṁ gatim || 9.32 || }
\cquote{O Pārtha, for those who take refuge in Me—even if they are of sinful birth, women, \emph{vaiśyas} (merchants), or \emph{śūdras} (labourers)—they too attain the supreme goal.}
\slcol{\Index{किं पुनर्ब्राह्मणाः पुण्या} भक्ता राजर्षयस्तथा । \\
अनित्यमसुखं लोकमिमं प्राप्य भजस्व माम् ॥ ९.३३ ॥ }
\translate{kiṁ punarbrāhmaṇāḥ puṇyā bhaktā rājarṣayastathā | \\
anityamasukhaṁ lōkamimaṁ prāpya bhajasva mām || 9.33 || }
\cquote{How much more, then, the holy \emph{brāhmaṇas} and devoted royal sages! Having come into this transient and joyless world, worship Me (with devotion).}
\slcol{\Index{मन्मना भव मद्भक्तो} मद्याजी मां नमस्कुरु । \\
मामेवैष्यसि युक्त्वैवमात्मानं मत्परायणः ॥ ९.३४ ॥ }
\translate{manmanā bhava madbhaktō madyājī māṁ namaskuru | \\
māmēvaiṣyasi yuktvaivamātmānaṁ matparāyaṇaḥ || 9.34 || }
\cquote{Fix your mind on Me, become My devotee, worship Me, bow down to Me. Thus, uniting the Self with Me, regarding Me as supreme, you shall surely come to Me.}

\newpage
\begin{mananam}{\mananamfont {Mananam Śloka - 34}}
\normalfont\mananamtext What changes can I make in my attitude to restore my divine status and connection with the ultimate Divine Being? How does bowing and serving a human authority differ from surrendering to the Divine? Could the Lord ever manipulate us for His own purposes as might an authoritative human? What kind of relationship exists between the Divine and me? Why is it considered the ultimate goal of life to re-establish or reclaim our true identity with the Divine?
\end{mananam}
\sreflect
\begin{inspiration}{\inspirationTitleFont{Inspiration}} 
\normalfont\mananamtext An immediate path to finding peace and adopting the right attitude is to inwardly surrender to life. Whether or not events go our way, people act as we wish, or our body and mind are in harmony, these factors are not always within our control. Yet, as life takes its course, we can cling to the one unchanging anchor of our lives—Whose love is unconditional and unchanged by life's circumstances.
\end{inspiration}
\newpage

\begin{center}
\devfont ॐ तत्सत्।  इति श्रिमद्भगवद्गीतासूपनिष्त्सु ब्रह्मविद्यायां योगशास्रे\\
श्रीकृष्णार्जुनसंवादे राजविद्याराजगुह्ययोगो नाम नवमोऽध्यायः ॥\\
\chapEndSlokaTranslate{ōṁ tatsaditi śrīmadbhagavadgītāsūpaniṣatsu brahmavidyāyāṁ yōgaśāstrē śrīkṛṣṇārjunasaṁvādē rājavidyārājaguhyayōgō nāma navamō'dhyāyaḥ||}
\end{center}